\documentclass[../../main.tex]{subfiles}
\begin{document}
\section{流形上的向量场}\label{section:ch08-流形上的向量场}

若 $M$ 是带边或不带边的光滑流形，则 $M$ 上的向量场是投影 $\pi\colon TM\to M$ 的一个截面。更具体地说，向量场是一个连续映射 $X\colon M\to TM$，通常写作 $p\mapsto X_p$，满足
\begin{align}
\pi\circ X=\operatorname{Id}_M,\label{align:ch08-01}
\end{align}
等价地，对每个 $p\in M$ 有 $X_p\in T_pM$。（我们将 $X$ 在 $p$ 处的值记作 $X_p$ 而非 $X(p)$，以与切丛元素的记号一致，同时避免与向量作用于函数的记号 $v(f)$ 混淆。）你应当以与 Euclid 空间中向量场相同的方式来想象 $M$ 上的向量场：即在 $M$ 的每一点附上一个箭头，取为与 $M$ 相切，且从一点到另一点连续变化。

\begin{figure}[ht]
\centering
\includegraphics[width=0.8\textwidth]{fig-08-1.png}
\caption{向量场}\label{figure:ch08-vector-field}
\end{figure}

我们主要关心光滑向量场，即作为从 $M$ 到 $TM$ 的映射是光滑的向量场，其中 $TM$ 赋予\cref{proposition:ch03-18}中描述的光滑流形结构。此外，出于某些目的，考虑从 $M$ 到 $TM$ 的、除可能不连续外即为向量场的映射是有用的。$M$ 上的\textbf{粗向量场}是满足\eqref{align:ch08-01}的一个（不必连续的）映射 $X\colon M\to TM$。与函数一样，若 $X$ 是 $M$ 上的向量场，则 $X$ 的\textbf{支撑集}定义为集合 $\{p\in M:X_p\neq0\}$ 的闭包。若向量场的支撑集是紧集，则称它是紧支撑的。

设 $M$ 是光滑 $n$ 流形（带边或不带边）。若 $X\colon M\to TM$ 是粗向量场，且 $\bigl(U,(x^i)\bigr)$ 是 $M$ 的任意光滑坐标卡，则可在坐标基向量下写出 $X$ 在任意点 $p\in U$ 处的值：
\begin{align}
X_p=X^i(p)\frac{\partial}{\partial x^i}\bigg|_p.\label{align:ch08-02}
\end{align}
这定义了 $n$ 个函数 $X^i\colon U\to\mathbb{R}$，称为 $X$ 在该坐标卡中的\textbf{分量函数}。

\begin{proposition}\label{proposition:ch08-01}[向量场的光滑性判据]
设 $M$ 是带边或不带边的光滑流形，$X\colon M\to TM$ 是粗向量场。若 $\bigl(U,(x^i)\bigr)$ 是 $M$ 上任意光滑坐标卡，则 $X$ 在 $U$ 上的限制是光滑的，当且仅当它在该坐标卡中的分量函数是光滑的。
\end{proposition}

\begin{proof}
设 $(x^i,v^i)$ 为 $\pi^{-1}(U)\subseteq TM$ 上与坐标卡 $(U,(x^i))$ 相伴的自然坐标。由自然坐标的定义，$X\colon M\to TM$ 在 $U$ 上的坐标表示为
\begin{align*}
\widehat{X}(x)=\bigl(x^1,\dots,x^n,X^1(x),\dots,X^n(x)\bigr),
\end{align*}
其中 $X^i$ 是 $X$ 在 $x^i$ 坐标下的第 $i$ 个分量函数。由此立即可知，$X$ 在 $U$ 上光滑等价于其分量函数光滑。
\end{proof}

\begin{example}\label{example:ch08-02}[坐标向量场]
若 $\bigl(U,(x^i)\bigr)$ 是 $M$ 上任意光滑坐标卡，则赋值
\begin{align*}
p\mapsto\frac{\partial}{\partial x^i}\bigg|_p
\end{align*}
确定了 $U$ 上的一个向量场，称为第 $i$ 个\textbf{坐标向量场}，记作 $\partial/\partial x^i$。它是光滑的，因为其分量函数为常数。
\end{example}

\begin{example}\label{example:ch08-03}[Euler 向量场]
$\mathbb{R}^n$ 上的向量场 $V$ 在 $x\in\mathbb{R}^n$ 处的值为
\begin{align*}
V_x=x^1\frac{\partial}{\partial x^1}\bigg|_x+\cdots+x^n\frac{\partial}{\partial x^n}\bigg|_x
\end{align*}
它是光滑的，因为其坐标函数是线性的。它在原点处为零，在其他各处沿径向向外。它称为 Euler 向量场，因为它出现在 Euler 齐次函数定理中（见\cref{exercise:ch08-02}）。
\end{example}

\begin{example}\label{example:ch08-04}[圆周上的角坐标向量场]
设 $\theta$ 是 $S^1$ 的适当开子集 $U\subseteq S^1$ 上任意角坐标（见\cref{exercise:ch01-08}），令 $\mathrm{d}/\mathrm{d}\theta$ 表示相应的坐标向量场。由于任何其他角坐标 $\widetilde{\theta}$ 在每一点的邻域内与 $\theta$ 相差一个加法常数，坐标向量场的变换法则 (3.11) 表明在它们的公共定义域上 $\mathrm{d}/\mathrm{d}\theta=\mathrm{d}/\mathrm{d}\widetilde{\theta}$。因此，$S^1$ 上存在一个全局定义的向量场，其关于任意角坐标的坐标表示为 $\mathrm{d}/\mathrm{d}\theta$。它是光滑向量场，因为其分量函数在任意此类坐标卡中为常数。我们将这个全局向量场记作 $\mathrm{d}/\mathrm{d}\theta$，尽管严格来说它不能同时被视为整个圆周上的坐标向量场。
\end{example}

\begin{example}\label{example:ch08-05}[环面上的角坐标向量场]
在 $n$ 维环面 $\mathbb{T}^n$ 上，为第 $i$ 个圆周因子选取角函数 $\theta_i$（$i=1,\dots,n$），便得到 $\mathbb{T}^n$ 的局部坐标 $(\theta^1,\dots,\theta^n)$。与前一例子类似的分析表明，坐标向量场 $\partial/\partial\theta_1,\dots,\partial/\partial\theta_n$ 在 $\mathbb{T}^n$ 上是光滑且全局定义的。
\end{example}

若 $U\subseteq M$ 是开集，则对每个 $p\in U$，$T_pU$ 自然等同于 $T_pM$（\cref{proposition:ch03-09}）这一事实使我们可以将 $TU$ 等同于 $TM$ 的开子集 $\pi^{-1}(U)$。因此，$U$ 上的向量场既可以视为从 $U$ 到 $TU$ 的映射，也可以视为从 $U$ 到 $TM$ 的映射，视何者更方便而定。若 $X$ 是 $M$ 上的向量场，则其限制 $X|_U$ 是 $U$ 上的向量场，且若 $X$ 光滑则它也光滑。

下一个引理是\cref{lemma:ch02-26}向向量场的推广，其证明方式大致相同。若 $M$ 是带边或不带边的光滑流形，$A\subseteq M$ 是任意子集，则沿 $A$ 的\textbf{向量场}是满足 $\pi\circ X=\operatorname{Id}_A$（换言之，对每个 $p\in A$ 有 $X_p\in T_pM$）的连续映射 $X\colon A\to TM$。若对每个 $p\in A$，存在 $p$ 在 $M$ 中的邻域 $V$ 以及 $V$ 上的光滑向量场 $\widetilde{X}$，使得在 $V\cap A$ 上 $\widetilde{X}$ 与 $X$ 一致，则称它为沿 $A$ 的光滑向量场。

\begin{lemma}\label{lemma:ch08-06}[向量场的延拓引理]
设 $M$ 是带边或不带边的光滑流形，$A\subseteq M$ 是闭子集。假设 $X$ 是沿 $A$ 的光滑向量场。给定任意包含 $A$ 的开子集 $U$，存在 $M$ 上的光滑全局向量场 $\widetilde{X}$，使得 $\widetilde{X}|_A=X$ 且 $\operatorname{supp}\widetilde{X}\subseteq U$。
\end{lemma}

\begin{proof}
见\cref{exercise:ch08-01}。
\end{proof}

作为重要的特殊情形，任意一点处的切向量都可以延拓为整个流形上的光滑向量场。

\begin{proposition}\label{proposition:ch08-07}
设 $M$ 是带边或不带边的光滑流形。给定 $p\in M$ 和 $v\in T_pM$，存在 $M$ 上的光滑全局向量场 $X$，使得 $X_p=v$。
\end{proposition}

\begin{proof}
赋值 $p\mapsto v$ 是上述定义中沿集合 $\{p\}$ 的向量场的一个例子。它是光滑的，因为它可以延拓为，例如，$p$ 的坐标邻域中的常系数向量场。因此，取 $A=\{p\}$ 和 $U=M$，由延拓引理即得本命题。
\end{proof}

若 $M$ 是带边或不带边的光滑流形，通常用记号 $\mathfrak{X}(M)$ 表示 $M$ 上所有光滑向量场的集合。它在逐点加法和数乘下构成向量空间：
\begin{align*}
(aX+bY)_p=aX_p+bY_p.
\end{align*}
该向量空间的零元素是零向量场，其在每个 $p\in M$ 处的值为 $0\in T_pM$。此外，光滑向量场可以乘以光滑实值函数：若 $f\in C^\infty(M)$ 且 $X\in\mathfrak{X}(M)$，则定义 $fX\colon M\to TM$ 为
\begin{align*}
(fX)_p=f(p)X_p.
\end{align*}
下一个命题表明这些运算产生光滑向量场。

\begin{proposition}\label{proposition:ch08-08}
设 $M$ 是带边或不带边的光滑流形。
\begin{enumerate}[(a)]
\item 若 $X$ 和 $Y$ 是 $M$ 上的光滑向量场，$f,g\in C^\infty(M)$，则 $fX+gY$ 是光滑向量场。
\item $\mathfrak{X}(M)$ 是 $C^\infty(M)$ 环上的模。
\end{enumerate}
\end{proposition}

\begin{exercise}\label{exercise:ch08-09}
证明\cref{proposition:ch08-08}。
\end{exercise}

例如，向量场 $X$ 的基表示\eqref{align:ch08-02}也可以写成向量场之间的等式，而非某点处向量之间的等式：
\begin{align*}
X=X^i\frac{\partial}{\partial x^i},
\end{align*}
其中 $X^i$ 是 $X$ 在给定坐标下的第 $i$ 个分量函数。

\textbf{局部标架与全局标架}

光滑坐标卡中的坐标向量场提供了一种表示向量场的便捷方式，因为它们的值在每一点构成切空间的一组基。然而，它们并非唯一选择。

设 $M$ 是带边或不带边的光滑 $n$ 流形。定义在某子集 $A\subseteq M$ 上的有序 $k$ 元向量场组 $(X_1,\dots,X_k)$ 称为\textbf{线性无关的}，若对每个 $p\in A$，$(X_1|_p,\dots,X_k|_p)$ 是 $T_pM$ 中线性无关的 $k$ 元组；称为\textbf{张成切丛}的，若对每个 $p\in A$，该 $k$ 元组张成 $T_pM$。$M$ 的\textbf{局部标架}是定义在开子集 $U\subseteq M$ 上的有序 $n$ 元向量场组 $(E_1,\dots,E_n)$，它线性无关且张成切丛；因而向量 $(E_1|_p,\dots,E_n|_p)$ 在每个 $p\in U$ 处构成 $T_pM$ 的一组基。若 $U=M$，则称为\textbf{全局标架}；若每个向量场 $E_i$ 都是光滑的，则称为\textbf{光滑标架}。我们常用简写记号 $(E_i)$ 表示标架 $(E_1,\dots,E_n)$。若 $M$ 的维数为 $n$，则要验证有序 $n$ 元向量场组 $(E_1,\dots,E_n)$ 是局部标架，只需验证它线性无关或张成切丛即可。

\begin{example}\label{example:ch08-10}[局部标架与全局标架]
\begin{enumerate}[(a)]
\item 标准坐标向量场构成 $\mathbb{R}^n$ 的光滑全局标架。

\item 若 $\bigl(U,(x^i)\bigr)$ 是光滑流形 $M$（可能带边）的任意光滑坐标卡，则坐标向量场在 $U$ 上构成光滑局部标架 $(\partial/\partial x^i)$，称为\textbf{坐标标架}。$M$ 的每一点都在此类局部标架的定义域内。

\item \cref{example:ch08-04}中定义的向量场 $\mathrm{d}/\mathrm{d}\theta$ 构成圆周的光滑全局标架。

\item \cref{example:ch08-04}中定义的 $n$ 维环面上的 $n$ 元向量场组 $(\partial/\partial\theta_1,\dots,\partial/\partial\theta_n)$ 是 $\mathbb{T}^n$ 的光滑全局标架。
\end{enumerate}
\end{example}

下一个命题表明局部标架很容易获得。

\begin{proposition}\label{proposition:ch08-11}[局部标架的补全]
设 $M$ 是带边或不带边的光滑 $n$ 流形。
\begin{enumerate}[(a)]
\item 若 $(X_1,\dots,X_k)$ 是开子集 $U\subseteq M$ 上线性无关的 $k$ 元光滑向量场组，$1\leqslant k<n$，则对每个 $p\in U$，存在 $p$ 的邻域 $V$ 中的光滑向量场 $X_{k+1},\dots,X_n$，使得 $(X_1,\dots,X_n)$ 在 $U\cap V$ 上是 $M$ 的光滑局部标架。

\item 若 $(v_1,\dots,v_k)$ 是某点 $p\in M$ 处 $T_pM$ 中线性无关的 $k$ 元向量组，$1\leqslant k\leqslant n$，则存在 $p$ 的邻域上的光滑局部标架 $(X_i)$，使得对 $i=1,\dots,k$ 有 $X_i|_p=v_i$。

\item 若 $(X_1,\dots,X_n)$ 是沿闭子集 $A\subseteq M$ 的线性无关的 $n$ 元光滑向量场组，则存在 $A$ 的某邻域上的光滑局部标架 $(\widetilde{X}_1,\dots,\widetilde{X}_n)$，使得对 $i=1,\dots,n$ 有 $\widetilde{X}_i|_A=X_i$。
\end{enumerate}
\end{proposition}

\begin{proof}
见\cref{exercise:ch08-05}。
\end{proof}

对于 $\mathbb{R}^n$ 的子集，有一种特殊类型的标架，在几何问题中通常比任意标架更有用。定义在某子集 $A\subseteq\mathbb{R}^n$ 上的 $k$ 元向量场组 $(E_1,\dots,E_k)$ 称为\textbf{正交规范的}，若对每个 $p\in A$，向量 $(E_1|_p,\dots,E_k|_p)$ 关于 Euclid 点积是正交规范的（其中我们以通常方式将 $T_p\mathbb{R}^n$ 等同于 $\mathbb{R}^n$）。由正交规范向量场组成的（局部或全局）标架称为\textbf{正交规范标架}。

\begin{example}\label{example:ch08-12}
标准坐标标架是 $\mathbb{R}^n$ 上的全局正交规范标架。作为一个不太显然的例子，考虑在 $\mathbb{R}^2\setminus\{0\}$ 上由
\begin{align}
E_1=\frac{x}{r}\frac{\partial}{\partial x}+\frac{y}{r}\frac{\partial}{\partial y},\qquad
E_2=-\frac{y}{r}\frac{\partial}{\partial x}+\frac{x}{r}\frac{\partial}{\partial y},\label{align:ch08-03}
\end{align}
定义的光滑向量场，其中 $r=\sqrt{x^2+y^2}$。直接计算表明 $(E_1,E_2)$ 是 $\mathbb{R}^2$ 在开子集 $\mathbb{R}^2\setminus\{0\}$ 上的正交规范标架。几何上，$E_1$ 和 $E_2$ 分别是与径向直线和以原点为圆心的圆相切的单位向量场。
\end{example}

下一个引理描述了一种构造正交规范标架的有用方法。

\begin{lemma}\label{lemma:ch08-13}[标架的 Gram--Schmidt 算法]
假设 $(X_j)$ 是 $T\mathbb{R}^n$ 在开子集 $U\subseteq\mathbb{R}^n$ 上的光滑局部标架。则 $U$ 上存在光滑正交规范标架 $(E_j)$，使得对每个 $j=1,\dots,n$ 和每个 $p\in U$，
\begin{align*}
\operatorname{span}\bigl(E_1|_p,\dots,E_j|_p\bigr)=\operatorname{span}\bigl(X_1|_p,\dots,X_j|_p\bigr).
\end{align*}
\end{lemma}

\begin{proof}
对每个 $p\in U$ 处的向量 $(X_j|_p)$ 应用 Gram--Schmidt 算法，得到由
\begin{align*}
E_j=\frac{X_j-\sum_{i=1}^{j-1}(X_j\cdot E_i)E_i}{\left|X_j-\sum_{i=1}^{j-1}(X_j\cdot E_i)E_i\right|}
\end{align*}
归纳给出的 $n$ 元粗向量场组 $(E_1,\dots,E_n)$。对每个 $j=1,\dots,n$ 和每个 $p\in U$，$X_j|_p\notin\operatorname{span}(E_1|_p,\dots,E_{j-1}|_p)$（它等于 $\operatorname{span}(X_1|_p,\dots,X_{j-1}|_p)$），故上述分母是 $U$ 上处处非零的光滑函数。因此，该公式将 $(E_j)$ 定义为 $U$ 上满足引理结论的光滑正交规范标架。
\end{proof}

虽然光滑局部标架很多，但全局标架并非如此。带边或不带边的光滑流形若容许光滑全局标架，则称为\textbf{可平行化的}。\cref{example:ch08-10}表明 $\mathbb{R}^n$、$S^1$ 和 $\mathbb{T}^n$ 都是可平行化的。\cref{exercise:ch08-06}和\cref{exercise:ch08-07}表明 $S^3$ 和 $S^7$ 是可平行化的。本章稍后将看到，所有 Lie 群都是可平行化的（见下文\cref{corollary:ch08-39}）。然而，尽管有这些例子，大多数光滑流形并不可平行化。（如我们将在第 10 章看到的，$M$ 的可平行化性与其切丛是否微分同胚于乘积 $M\times\mathbb{R}^n$ 的问题密切相关。）

不可平行化流形最简单的例子是 $S^2$，但这一事实的证明须等到我们发展更多工具之后（见\cref{exercise:ch16-06}）。事实上，利用代数拓扑中更高级的方法，Raoul Bott 和 John Milnor [MB58] 以及 Michel Kervaire [Ker58] 于 1958 年独立证明了 $S^1$、$S^3$ 和 $S^7$ 是仅有的可平行化球面。因此，它们是仅有的可能容许 Lie 群结构的正维球面。前两个确实容许（见\cref{example:ch07-10}(a)和\cref{exercise:ch07-16}），但结果表明 $S^7$ 没有 Lie 群结构（见 [Bre93, p. 301]）。

\textbf{作为 $C^\infty(M)$ 导子的向量场}

向量场的一个基本性质是它们在光滑实值函数空间上定义算子。若 $X\in\mathfrak{X}(M)$，$f$ 是定义在开子集 $U\subseteq M$ 上的光滑实值函数，则得到新函数 $Xf\colon U\to\mathbb{R}$，定义为
\begin{align*}
(Xf)(p)=X_pf.
\end{align*}
（注意不要混淆记号 $fX$ 和 $Xf$：前者是将 $X$ 乘以 $f$ 得到的 $U$ 上的光滑向量场，后者是将向量场 $X$ 作用于光滑函数 $f$ 得到的 $U$ 上的实值函数。）由于切向量对函数的作用由函数在任意小邻域内的值决定，由此可知 $Xf$ 是局部确定的。特别地，对任意开子集 $V\subseteq U$，
\begin{align}
(Xf)|_V=X(f|_V).\label{align:ch08-04}
\end{align}

这一构造给出了向量场的另一个有用的光滑性判据。

\begin{proposition}\label{proposition:ch08-14}
设 $M$ 是带边或不带边的光滑流形，$X\colon M\to TM$ 是粗向量场。下列条件等价：
\begin{enumerate}[(a)]
\item $X$ 是光滑的。
\item 对每个 $f\in C^\infty(M)$，函数 $Xf$ 在 $M$ 上光滑。
\item 对每个开子集 $U\subseteq M$ 和每个 $f\in C^\infty(U)$，函数 $Xf$ 在 $U$ 上光滑。
\end{enumerate}
\end{proposition}

\begin{proof}
我们证明 (a) $\Rightarrow$ (b) $\Rightarrow$ (c) $\Rightarrow$ (a)。

为证明 (a) $\Rightarrow$ (b)，假设 $X$ 光滑，$f\in C^\infty(M)$。对任意 $p\in M$，可在 $p$ 的邻域 $U$ 中选取光滑坐标 $(x^i)$。则对 $x\in U$，可写
\begin{align*}
Xf(x)=\left(X^i(x)\frac{\partial}{\partial x^i}\bigg|_x\right)f=X^i(x)\frac{\partial f}{\partial x^i}(x).
\end{align*}
由于由\cref{proposition:ch08-01}可知分量函数 $X^i$ 在 $U$ 上光滑，由此可知 $Xf$ 在 $U$ 上光滑。由于在每一点的邻域内都如此，$Xf$ 在 $M$ 上光滑。

为证明 (b) $\Rightarrow$ (c)，假设 $U\subseteq M$ 是开集，$f\in C^\infty(U)$。对任意 $p\in U$，取光滑鼓函数 $\psi$，它在 $p$ 的邻域内等于 $1$ 且支撑在 $U$ 中，并令 $\widetilde{f}=\psi f$，在 $M\setminus\operatorname{supp}\psi$ 上延拓为零。则由假设 $X\widetilde{f}$ 光滑，且由\eqref{align:ch08-04}它在 $p$ 的邻域内等于 $Xf$。这表明 $Xf$ 在 $U$ 的每一点附近光滑。

最后，为证明 (c) $\Rightarrow$ (a)，假设每当 $f$ 在 $M$ 的开子集上光滑时 $Xf$ 光滑。若 $(x^i)$ 是 $U\subseteq M$ 上任意光滑局部坐标，则可将每个坐标 $x^i$ 视为 $U$ 上的光滑函数。将 $X$ 作用于这些函数之一，得到
\begin{align*}
Xx^i=X^j\frac{\partial}{\partial x^j}(x^i)=X^i.
\end{align*}
由于由假设 $Xx^i$ 光滑，由此可知 $X$ 的分量函数光滑，故 $X$ 光滑。
\end{proof}

上述命题的一个推论是，光滑向量场 $X\in\mathfrak{X}(M)$ 通过 $f\mapsto Xf$ 定义了从 $C^\infty(M)$ 到自身的映射。该映射显然是 $\mathbb{R}$-线性的。此外，切向量的乘积法则 (3.4) 转化为向量场的如下乘积法则：
\begin{align}
X(fg)=fXg+gXf,\label{align:ch08-05}
\end{align}
你可以通过在任意点 $p\in M$ 处计算两边来验证。一般地，映射 $X\colon C^\infty(M)\to C^\infty(M)$ 称为\textbf{导子}（不同于第 3 章中定义的 $p$ 处的导子），若它是 $\mathbb{R}$-线性的且对所有 $f,g\in C^\infty(M)$ 满足\eqref{align:ch08-05}。

下一个命题表明 $C^\infty(M)$ 的导子可以与光滑向量场等同。

\begin{proposition}\label{proposition:ch08-15}
设 $M$ 是带边或不带边的光滑流形。映射 $D\colon C^\infty(M)\to C^\infty(M)$ 是导子，当且仅当它形如 $Df=Xf$，其中 $X\in\mathfrak{X}(M)$ 是某个光滑向量场。
\end{proposition}

\begin{proof}
我们刚刚证明了每个光滑向量场诱导一个导子。反过来，假设 $D\colon C^\infty(M)\to C^\infty(M)$ 是导子。我们需要构造一个向量场 $X$，使得对所有 $f$ 有 $Df=Xf$。由上面的讨论，显然若存在这样的向量场，其在 $p\in M$ 处的值必须是 $p$ 处的导子，其对任意光滑实值函数 $f$ 的作用由
\begin{align*}
X_pf=(Df)(p)
\end{align*}
给出。$D$ 的线性性保证该表达式线性地依赖于 $f$，而 $D$ 是导子这一事实给出切向量的乘积法则 (3.4)。因此，如此定义的映射 $X_p\colon C^\infty(M)\to\mathbb{R}$ 确实是切向量，即 $C^\infty(M)$ 在 $p$ 处的导子。这将 $X$ 定义为粗向量场。因为每当 $f\in C^\infty(M)$ 时 $Xf=Df$ 光滑，故由\cref{proposition:ch08-14}可知该向量场光滑。
\end{proof}

由于这一结果，我们有时将 $M$ 上的光滑向量场与 $C^\infty(M)$ 的导子等同，对向量场（视为从 $M$ 到 $TM$ 的光滑映射）和导子（视为从 $C^\infty(M)$ 到自身的线性映射）使用同一字母。

\end{document}